% ===========================================================================
%  Devoir surveillé nº 1 — fin de la partie Algèbre
% ===========================================================================
\chapter*{Devoir surveillé nº\,1 — Algèbre}
\markboth{Devoir surveillé nº\,1}{Devoir surveillé nº\,1}
\addcontentsline{toc}{chapter}{Devoir surveillé nº\,1 — Algèbre}

\begin{devoirsurveille}[portée : chapitres 1 et 2]
\textbf{Durée : 2 h 15.} Les deux exercices et le problème sont obligatoires.
Calculatrice scientifique non programmable autorisée. Le barème est donné
pour les deux options : \textup{(A1 ; A2)}.

\vspace{0.7em}
\noindent\textbf{\sffamily EXERCICE 1} \hfill \textup{(05 points)}

\begin{enumerate}[label=\textbf{\arabic*.}, leftmargin=*, itemsep=0.25em]
  \item Résoudre dans $\R$ l'équation $2x^{2}-7x+3 = 0$. \bareme{1 ; 1}
  \item En déduire une factorisation de $2x^{2}-7x+3$. \bareme{0,5 ; 0,5}
  \item Résoudre dans $\R$ l'inéquation $2x^{2}-7x+3 \leq 0$. \bareme{1 ; 1}
  \item Résoudre dans $\R$ l'équation $2x^{4}-7x^{2}+3 = 0$.
        \bareme{1,5 ; 1,5}
  \item Pour quelles valeurs de $x$ l'expression
        $\sqrt{2x^{2}-7x+3}$ est-elle définie ? \bareme{1 ; 1}
\end{enumerate}

\vspace{0.7em}
\noindent\textbf{\sffamily EXERCICE 2} \hfill \textup{(05 points)}

On considère le système
\[
  (S) \quad \begin{cases} 3x-2y = 4\\ 5x+y = 19 \end{cases}
\]

\begin{enumerate}[label=\textbf{\arabic*.}, leftmargin=*, itemsep=0.25em]
  \item Calculer le déterminant $D$ de $(S)$ et dire ce qu'on en conclut.
        \bareme{1 ; 1}
  \item Résoudre $(S)$ par la méthode de Cramer. \bareme{2 ; 2}
  \item Vérifier le couple trouvé dans les deux équations. \bareme{0,5 ; 0,5}
  \item Déterminer le réel $m$ pour que le système
        $\begin{cases} 3x-2y = 4\\ mx+y = 19\end{cases}$
        n'admette pas de solution unique. \bareme{1,5 ; 1,5}
\end{enumerate}

\vspace{0.7em}
\noindent\textbf{\sffamily PROBLÈME} \hfill \textup{(10 points)}

Une association scolaire d'Antsirabe organise une kermesse. Elle vend trois
sortes de billets : billet enfant, billet élève et billet adulte. On note
$x$, $y$ et $z$ le nombre de billets vendus dans chacune de ces trois
catégories.

\medskip
\noindent\textbf{Partie A — mise en équations}

\begin{enumerate}[label=\textbf{\arabic*.}, leftmargin=*, itemsep=0.25em]
  \item L'association a vendu $500$ billets en tout. Traduire cette phrase
        par une équation. \bareme{0,5 ; 0,5}
  \item Elle a vendu deux fois plus de billets élève que de billets enfant.
        Traduire cette phrase par une équation. \bareme{0,5 ; 0,5}
  \item Le billet enfant coûte $1\,000$ ariary, le billet élève $2\,000$
        ariary et le billet adulte $5\,000$ ariary ; la recette totale est de
        $1\,450\,000$ ariary. Traduire cette phrase par une équation, puis
        la simplifier en divisant par $1\,000$. \bareme{1 ; 1}
  \item Écrire le système $(T)$ ainsi obtenu, ordonné en $x$, $y$, $z$.
        \bareme{0,5 ; 0,5}
\end{enumerate}

\medskip
\noindent\textbf{Partie B — résolution}

\begin{enumerate}[label=\textbf{\arabic*.}, leftmargin=*, itemsep=0.25em, start=5]
  \item Calculer le déterminant de $(T)$ par la règle de Sarrus.
        \bareme{2 ; 1,5}
  \item Que peut-on en conclure quant au nombre de solutions ?
        \bareme{0,5 ; 0,5}
  \item Résoudre $(T)$ et donner le nombre de billets de chaque sorte.
        \bareme{2,5 ; 2,5}
  \item Vérifier la recette totale à partir du triplet trouvé.
        \bareme{0,5 ; 0,5}
\end{enumerate}

\medskip
\noindent\textbf{Partie C — une contrainte de plus}

\begin{enumerate}[label=\textbf{\arabic*.}, leftmargin=*, itemsep=0.25em, start=9]
  \item L'an prochain, l'association veut vendre $n$ billets adulte de plus
        et autant de billets enfant en moins, le nombre de billets élève
        restant inchangé. Exprimer la nouvelle recette, en milliers
        d'ariary, en fonction de $n$. \bareme{1 ; 1}
  \item Pour quelles valeurs entières de $n$ la recette dépassera-t-elle
        $1\,600\,000$ ariary, sachant qu'on ne peut pas vendre un nombre
        négatif de billets enfant ? \bareme{1,5 ; 1,5}
\end{enumerate}
\end{devoirsurveille}

\begin{corrige}[devoir surveillé nº 1 — exercice 1]
\textbf{1.} $\Delta = 49-24 = 25$, $\sqrt\Delta = 5$ :
$x_{1} = \frac{7-5}{4} = \frac12$ et $x_{2} = \frac{7+5}{4} = 3$.

\textbf{2.} $2x^{2}-7x+3 = 2\left(x-\frac12\right)(x-3) = (2x-1)(x-3)$.

\textbf{3.} Le coefficient dominant est positif : le trinôme est négatif
entre les racines. $\mathcal S = \intervalle{\frac12}{3}$.

\textbf{4.} En posant $X = x^{2}$ : $X = \frac12$ ou $X = 3$, toutes deux
positives, d'où
$x \in \left\{-\sqrt3\,;\,-\frac{\sqrt2}{2}\,;\,\frac{\sqrt2}{2}\,;\,
\sqrt3\right\}$.

\textbf{5.} La racine carrée impose $2x^{2}-7x+3 \geq 0$, c'est-à-dire
$x \in \intervalleof{-\infty}{\frac12}\cup\intervallefo{3}{+\infty}$.
\end{corrige}

\begin{corrige}[devoir surveillé nº 1 — exercice 2]
\textbf{1.} $D = 3\times1-5\times(-2) = 13 \neq 0$ : le système admet un
unique couple solution.

\textbf{2.} $D_{x} = 4\times1-19\times(-2) = 42$ et
$D_{y} = 3\times19-5\times4 = 37$. Donc $x = \frac{42}{13}$ et
$y = \frac{37}{13}$.

\textbf{3.} $3\times\frac{42}{13}-2\times\frac{37}{13}
= \frac{126-74}{13} = \frac{52}{13} = 4$ \checkmark ; et
$5\times\frac{42}{13}+\frac{37}{13} = \frac{210+37}{13} = \frac{247}{13}
= 19$ \checkmark.

\textbf{4.} $D = 3+2m = 0$ pour $m = -\frac32$.
\end{corrige}

\begin{corrige}[devoir surveillé nº 1 — problème]
\textbf{Partie A.}
\textbf{1.} $x+y+z = 500$.
\textbf{2.} $y = 2x$, soit $2x-y = 0$.
\textbf{3.} $1\,000x+2\,000y+5\,000z = 1\,450\,000$, soit
$x+2y+5z = 1\,450$.
\textbf{4.}
$(T) : \begin{cases} x+y+z = 500\\ 2x-y+0\,z = 0\\ x+2y+5z = 1\,450
\end{cases}$

\textbf{Partie B.}
\textbf{5.} Par Sarrus, sur les coefficients
$\begin{vmatrix} 1 & 1 & 1\\ 2 & -1 & 0\\ 1 & 2 & 5\end{vmatrix}$ :
les produits descendants donnent $(1)(-1)(5)+(1)(0)(1)+(1)(2)(2) = -5+0+4
= -1$ ; les produits montants donnent
$(1)(-1)(1)+(1)(2)(0)+(1)(5)(2) = -1+0+10 = 9$. D'où $D = -1-9 = -10$.

\textbf{6.} $D \neq 0$ : le système admet un unique triplet solution.

\textbf{7.} La deuxième équation donne $y = 2x$. En reportant :
$3x+z = 500$ et $5x+5z = 1\,450$, soit $x+z = 290$. En soustrayant,
$2x = 210$, donc $x = 105$, $y = 210$ et $z = 185$. L'association a vendu
$105$ billets enfant, $210$ billets élève et $185$ billets adulte.

\textbf{8.} $105+2\times210+5\times185 = 105+420+925 = 1\,450$, soit
$1\,450\,000$ ariary. \checkmark

\textbf{Partie C.}
\textbf{9.} Les effectifs deviennent $105-n$, $210$ et $185+n$. La recette,
en milliers d'ariary, vaut
\[
  R(n) = (105-n)+2\times210+5(185+n) = 1\,450+4n .
\]
\textbf{10.} $R(n) > 1\,600$ équivaut à $4n > 150$, soit $n > 37{,}5$, donc
$n \geq 38$. Par ailleurs $105-n \geq 0$ impose $n \leq 105$. Les valeurs
entières cherchées sont donc celles de $38$ à $105$.
\end{corrige}

\begin{notepourlaclasse}
Le problème est calibré pour une classe qui a fait les deux chapitres. Sa
partie C est celle qui départage : elle demande de reconstruire une recette
à partir d'effectifs modifiés, ce que peu d'élèves font spontanément sans un
tableau. Suggérez-leur d'en tracer un au brouillon — trois lignes, trois
colonnes — avant de calculer.

Barème : les deux options composent sur le même sujet ; la colonne A2 est
plus resserrée sur la question 5, qui est mécanique, pour laisser du temps à
la partie C.

Si le devoir doit tenir en une heure, gardez l'exercice 1 en entier et les
parties A et B du problème : cela fait exactement dix points et couvre les
deux chapitres.
\end{notepourlaclasse}
